\documentclass[11pt,a4paper]{article}
\usepackage[margin=25mm,headheight=14pt]{geometry}
\usepackage[T1]{fontenc}
\usepackage[utf8]{inputenc}
\usepackage{lmodern,amsmath,amssymb}
\usepackage{xcolor,graphicx,booktabs,longtable,array}
\usepackage{microtype,enumitem,fancyhdr,titlesec}
\usepackage{hyperref}
\definecolor{accent}{HTML}{163B55}
\hypersetup{colorlinks=true,linkcolor=accent,urlcolor=accent,
 pdftitle={A counterexample to an additive intersection-index eigenvalue bound},pdfauthor={Matthew J. Colbrook},
 pdfsubject={AIM 1022: mathematical result and proof}}
\urlstyle{same}
\Urlmuskip=0mu plus 1mu
\setlength{\parindent}{0pt}
\setlength{\parskip}{0.35em}
\setlength{\emergencystretch}{3em}
\setlist{itemsep=0.25em,topsep=0.3em}
\providecommand{\tightlist}{\setlength{\itemsep}{0pt}\setlength{\parskip}{0pt}}
\titleformat{\section}{\normalfont\large\bfseries\color{accent}}{\thesection}{0.7em}{}
\titleformat{\subsection}{\normalfont\normalsize\bfseries}{\thesubsection}{0.6em}{}
\pagestyle{fancy}
\fancyhf{}
\fancyhead[L]{\small\sffamily AIM 1022}
\fancyhead[R]{\small\sffamily Matthew J. Colbrook}
\fancyfoot[C]{\thepage}
\renewcommand{\headrulewidth}{0.3pt}
\setcounter{secnumdepth}{0}
\allowdisplaybreaks[2]
\title{\textbf{A counterexample to an additive intersection-index eigenvalue bound}}
\author{Matthew J. Colbrook\\[3pt]
{\small Department of Applied Mathematics and Theoretical Physics}\\
{\small University of Cambridge, United Kingdom}\\
{\small\href{mailto:mjc249@maths.cam.ac.uk}{\texttt{mjc249@maths.cam.ac.uk}}}}
\date{3 November 2026}
\begin{document}
\maketitle
\begin{abstract}
We construct embedded hypersurfaces that contradict the additive intersection-index spectral estimate in AIM 1022. The argument combines a base with large normalized first eigenvalue, tapered connecting handles, a uniform intersection count and Neumann spectral counting.
\end{abstract}
\noindent\textbf{Repository target.}
This manuscript concerns AIM problem 1022 at revision
\texttt{3fa50884d129}. Its scope is stated below.
\medskip

\textbf{Theorem.} For every integer \(n\ge 2\), there do not exist a
finite constant \(A_n>0\) and a continuous function
\(F_n:[0,\infty)\to\mathbb R\) such that every smooth compact connected
embedded hypersurface \(\Sigma ^n\) in \(\mathbb R^{n+1}\), with finite
intersection index, satisfies

\[
 \lambda_k(\Sigma)\operatorname{Vol}_n(\Sigma)^{2/n}
 \le F_n(i(\Sigma))+A_n k^{2/n}\qquad(k\ge1).
\]

The additive question is posed by P\'etiard [5, p.~3].
The proof replicates a fixed hypersurface with a large normalized first
eigenvalue, joining its copies by long handles whose Neumann spectra
contain at most two low modes each. The copies and handles can be placed
so that the intersection index stays bounded as their number increases.
The fixed base may have a large index; it is independence from the
replication number that matters.

The only substantial external inputs are the existence of embedded bases
with unbounded normalized first eigenvalue {[}1, Theorem 1.4{]}, and
uniform finiteness for definable families in the o-minimal structure
\(\mathbb R_{\mathrm{an},\exp}\) {[}2,3{]}. All handle, placement,
smoothing, and spectral-counting details are supplied below. This is a
mathematical proof developed in this research session, not a claim of
external peer review or priority.

\hypertarget{1022-bases-with-large-normalized-first-eigenvalue-and-finite-definable-complexity}{%
\section{0. Bases with large normalized first eigenvalue and finite
definable
complexity}\label{1022-bases-with-large-normalized-first-eigenvalue-and-finite-definable-complexity}}

Colbois--Dryden--El Soufi {[}1, Theorem 1.4{]} prove that, for every
\(n\ge 2\),

\[
 \sup_B \lambda_1(B)\operatorname{Vol}_n(B)^{2/n}=\infty,
 \tag{0.1}
\]

where B ranges over smooth compact connected embedded hypersurfaces of
\(\mathbb R^{n+1}\). In dimension two the topology may vary; in higher
dimensions the result is stronger, but that strengthening is unnecessary
here.

We can require each chosen base to be a regular compact component of a
polynomial zero set. Here are details of the approximation. A smooth
closed embedded hypersurface has a compact tubular neighborhood on which
its signed distance rho is smooth, with nonzero normal derivative.
Extend rho smoothly to a surrounding compact box. Approximate this
extension in \(C^2\) on the box by polynomials. For a sufficiently close
polynomial p, its zero set in the tubular neighborhood is a unique
\(C^1\)-small normal graph over B, with nonvanishing gradient. The
polynomial has fixed nonzero signs on the two boundary sheets of a
sufficiently small tube, so this zero component cannot leave the tube or
join any other zero component. It is therefore a smooth compact
connected component of the real algebraic set \({p=0}\).

Such a component is semialgebraic. Other components of the polynomial
zero set are irrelevant: the base is only the selected compact
component. The induced metrics converge uniformly under the \(C^1\)
normal-graph convergence. The variational characterization of
eigenvalues and the volume formula then give convergence of
\(\lambda _1\) and volume. Hence (0.1) still holds for these
semialgebraic bases.

After a Euclidean dilation, normalize

\[
 \lambda_1(B)=1.
\]

The base volume \(V=\operatorname{Vol} _n(B)\) can consequently be
arbitrarily large, because this rescaling changes the volume to the
\(n/2\) power of the original normalized first eigenvalue. Fix one such
base for the rest of the construction. Its definable complexity, and all
constants depending on it, are allowed to be large but remain fixed as
the number N of copies tends to infinity.

\hypertarget{1022-a.-a-straight-handle-with-at-most-two-low-neumann-eigenvalues}{%
\section{A. A straight handle with at most two low Neumann
eigenvalues}\label{1022-a.-a-straight-handle-with-at-most-two-low-neumann-eigenvalues}}

Fix \(a>0\), length \(L>0\), and radius parameter \(\varepsilon >0\).
Put

\[
 r(s)=\epsilon\frac{\cosh(a(s-L/2))}{\cosh(aL/2)},\qquad0\le s\le L.
\]

The \(n\)-dimensional tube of revolution in \(\mathbb R^{n+1}\) has
metric

\[
 g=q(s)^2 ds^2+r(s)^2g_{S^{n-1}},\qquad
 q(s)=\sqrt{1+r'(s)^2}\le\sqrt{1+a^2\epsilon^2}=:Q.
\]

With Neumann conditions at the two end cross-sections, separate
spherical harmonics. In angular degree zero the Rayleigh quotient is

\[
 \frac{\int_0^L r^{n-1}q^{-1}|f'|^2\,ds}
      {\int_0^L qr^{n-1}|f|^2\,ds}
 \ge Q^{-2}
 \frac{\int_0^L r^{n-1}|f'|^2\,ds}
      {\int_0^L r^{n-1}|f|^2\,ds}.
 \tag{A1}
\]

Write \(p=(n-1)/2\). On the codimension-two subspace \(H^1_0(0,L)\),
putting \(\psi =r^p f\) and integrating by parts gives

\[
 \int_0^Lr^{2p}|f'|^2
 =\int_0^L\left(|\psi'|^2+
       \frac{(r^p)''}{r^p}|\psi|^2\right),
\]

where

\[
 \frac{(r^p)''}{r^p}
  =a^2\left[p\operatorname{sech}^2(a(s-L/2))
              +p^2\tanh^2(a(s-L/2))\right]
  \ge a^2\min(p,p^2).
 \tag{A2}
\]

The min--max principle therefore says that the third angular-zero
Neumann eigenvalue (that is, the eigenvalue with index 2 when counting
from zero) is at least

\[
 \Gamma_\epsilon=
 \frac{a^2\min(p,p^2)}{1+a^2\epsilon^2}.
 \tag{A3}
\]

Equivalently, that sector has at most two eigenvalues strictly below
\(\Gamma _\varepsilon\). Every nonconstant spherical mode has Rayleigh
quotient at least

\[
 \frac{n-1}{\epsilon^2},
 \tag{A4}
\]

because its angular eigenvalue is at least \(n-1\) and
\(r\le \varepsilon\). Thus, after fixing a sufficiently large a and then
taking epsilon small, \textbf{the entire handle has at most two
eigenvalues below any prescribed fixed threshold}, independently of its
length L.

Its volume also satisfies a length-independent estimate. Since

\[
 \frac{\cosh(a(s-L/2))}{\cosh(aL/2)}
 \le e^{-as}+e^{-a(L-s)},
\]

and \((x+y)^{n-1}\le 2^{n-2}(x^{n-1}+y^{n-1})\),

\[
 \operatorname{Vol}_n(\text{handle})
 \le Q|S^{n-1}|\frac{2^{n-1}}{(n-1)a}\epsilon^{n-1}.
 \tag{A5}
\]

\hypertarget{1022-b.-each-straight-handle-meets-an-affine-line-at-most-four-times}{%
\section{B. Each straight handle meets an affine line at most four
times}\label{1022-b.-each-straight-handle-meets-an-affine-line-at-most-four-times}}

If a line has nonconstant axial coordinate s, its transverse coordinate
is \(A+Bs\). Its intersections with the tube satisfy

\[
 |A+Bs|^2-r(s)^2=0.
\]

The fourth derivative of the left side is strictly negative, since
\((r^2)''''>0\). Five distinct zeros would contradict Rolle's theorem.
If the axial coordinate is constant, the line meets one spherical
cross-section in at most two points. Hence the handle contributes at
most four transverse intersections to any affine line.

\hypertarget{1022-c.-a-robust-bound-of-eight-handles-intersected-by-one-line}{%
\section{C. A robust bound of eight handles intersected by one
line}\label{1022-c.-a-robust-bound-of-eight-handles-intersected-by-one-line}}

Let \(m=n+1\ge 3\). Choose a generic polygonal chain with N vertices in
\(\mathbb R^m\), with all increments in a prescribed open cone around
\(+e_m\). The vertices may additionally be restricted to an open set
with consecutive edge lengths between 1 and 2.

For five \textbf{pairwise vertex-disjoint} edges, the set of vertex
configurations admitting a common affine-line transversal has positive
codimension. Indeed affine lines have dimension \(2m-2\); for a fixed
line, one segment's endpoints allowing an intersection have dimension at
most \(m+2\) rather than \(2m\). Because the five endpoint pairs are
disjoint, these conditions are independent. The total incidence-set
dimension is at most

\[
 mN+(2m-2)-5(m-2)<mN
\]

for every \(m\ge 3\). This parametrization also includes endpoint
intersections; those strata have no larger dimension. There are only
finitely many choices of five edges, so generic vertices avoid all such
configurations.

Every collection of nine edges of a chain contains a matching of five
pairwise vertex-disjoint edges. Therefore no affine line meets nine
edges of this generic chain. A line meets \textbf{at most eight} core
edges. This slightly weaker constant avoids degeneracies involving
adjacent edges and is fully sufficient.

We also choose the vertices so no three are collinear and nonadjacent
core edges do not intersect. These are additional finite genericity
conditions.

For the fixed finite chain, all these nonintersection conditions have a
positive tolerance. A line meeting a fixed bounded neighborhood of the
chain belongs to a compact line-parameter space after restricting its
offset. Thus, if arbitrarily small tubular neighborhoods of nine chosen
edges had a common transversal, compactness would give a transversal of
their cores. Taking the minimum over the finitely many subsets gives a
positive tube-radius tolerance. The same compactness argument supplies a
positive radius tolerance for no line meeting three vertex balls.

Global dilation of the prototype chain increases these margins
arbitrarily. Consequently its vertices can be replaced by fixed-size
base hypersurfaces, and its edges by arbitrarily long thin handles,
while retaining the bounds of two base balls and eight handles per line.

\hypertarget{1022-d.-matching-the-actual-ports-to-the-prototype-chain}{%
\section{D. Matching the actual ports to the prototype
chain}\label{1022-d.-matching-the-actual-ports-to-the-prototype-chain}}

Let the fixed base hypersurface be a regular compact component of a
polynomial zero set. Choose a direction v such that both v and -v are
regular values of its Gauss map and the height in direction v has a
unique global maximum and a unique global minimum. This is possible:
Sard's theorem gives the first condition almost everywhere, while
almost-everywhere differentiability of the support function of the
convex hull gives uniqueness of the two support points almost
everywhere. At these extrema, nondegeneracy implies strict convexity of
the outward second fundamental form.

The normal maps on small neighborhoods of these two points therefore
have analytic semialgebraic inverses

\[
 p_+(d),\qquad p_-(-d),
\]

for d in a sufficiently small open cone around v. The rays leaving these
ports in their outward normals avoid the rest of the fixed compact base,
after narrowing the cone if necessary.

Let the prototype chain have increments \(L_j d_j\). For a dilation
factor R, choose actual translations recursively by

\[
 t_{j+1}=t_j+p_+(d_j)+R L_jd_j-p_-(-d_j).
 \tag{D1}
\]

The outgoing port on copy j and incoming port on copy j+1 are then
exactly joined by a straight segment of length \(R L_j\) in direction
d\_j. The actual translations and segments differ from those of the
dilated prototype chain by at most a constant times N, because the port
offsets are uniformly bounded. For each fixed N, choose R so large that
these bounded errors are smaller than the dilated nonintersection
margins. This preserves the no-three-base-balls and no-nine-core-edges
conclusions. All copies use the same base metric and orientation.

The smooth collars below occupy a fixed-radius base ball, and the
straight tube begins a distance epsilon beyond each port. Thus its
length is \(R L_j-2\varepsilon\) rather than \(R L_j\); this harmless
change is already covered by the estimates in A.

\hypertarget{1022-e.-explicit-definable-collars-without-integrating-a-cutoff}{%
\section{E. Explicit definable collars, without integrating a
cutoff}\label{1022-e.-explicit-definable-collars-without-integrating-a-cutoff}}

This construction supplies the local family needed for the
intersection-index argument. In coordinates centered at a port, with
outward normal the z-axis, write the base locally as an analytic
semialgebraic graph \(z=g(y)\), where \(g(0)=0\) and \(Dg(0)=0\). Its
interior lies below the graph.

Choose a fixed nondecreasing \(C^\infty\) definable cutoff chi on
\([0,1]\), equal to zero on \([0,1/4]\) and one on \([3/4,1]\). For
example it can be built from the flat function \(\exp (-1/t)\) for
\(t>0\), extended by zero for \(t\le 0\), using products and quotients;
no integration is needed. Put \(\eta (t)=t \chi (t)\). Then eta is zero
near zero, equals t near one, and has positive derivative when it is not
zero.

A fixed corner profile is

\[
 \rho_0(t)=\epsilon[1+\eta(1-t)],\qquad
 z_0(t)=\epsilon\eta(t),\qquad 0\le t\le1.
\]

It is planar near t=0 and cylindrical near t=1, with no zero tangent
vector. Let \(\chi_{\mathrm{in}}\) be another such cutoff, zero on
\([0,1/2]\) and one on \([3/4,1]\). Let r be the cosh tube profile,
analytically continued a short distance to negative s. Define

\[
 \rho(t)=\rho_0(t)+\chi_{\rm in}(t)[r(z_0(t)-\epsilon)-\epsilon].
\]

Choose \(\chi_{\mathrm{base}}\) equal to one near zero and zero on
\([1/2,1]\). The collar is parametrized by

\[
 (t,\theta)\longmapsto
 \left(\rho(t)\theta,
       z_0(t)+\chi_{\rm base}(t)g(\rho(t)\theta)\right),
 \qquad\theta\in S^{n-1}.
 \tag{E1}
\]

Near t=0 this agrees exactly with the base graph on an annulus. Near t=1
it agrees exactly with the continued tube, with axial coordinate
\(s=z_0(t)-\varepsilon\); thus all derivatives match at both ends.
Remove the base disk of radius \(2\varepsilon\) and retain the collar up
to t=1, where the straight tube begins.

Because \(g(y)=O(|y|^2)\) and the cosh correction on a segment of length
\(O(\varepsilon )\) is \(O(a \varepsilon ^2)\), after rescaling by
epsilon this is a small \(C^1\) perturbation of the fixed corner
profile. It is embedded and has a uniformly controlled scaled metric for
sufficiently small epsilon. These bounds are uniform in the tube length:
on the short continued segment,
\(|r^{(j)}(s)|\le a^j\varepsilon e^{a\varepsilon}\). The fixed corner
profile is a compact embedded annular cylinder with nonzero tangent;
sufficiently small C1 perturbations preserve its embedding. One may also
choose chi\_base nonincreasing, in which case strict convexity of the
port graph makes its axial coordinate monotone through the outer
transition. The same construction at the other endpoint uses its outward
normal, which reverses the tube coordinate; the cosh profile is
symmetric about its midpoint.

All functions and parametrizations above are definable in
\(\mathbb R_{\mathrm{an},\exp}\). The base is fixed, its local graph and
inverse Gauss maps are semialgebraic, coordinate rotations can be chosen
semialgebraically in the normal cone, and the cutoff and cosh terms use
only exponentials and algebraic operations. This is a fixed
finite-parameter family: parameters are the two port directions,
epsilon, a, and the two incident handle lengths \(L_-\) and \(L_+\)
(with a fixed after the spectral threshold is chosen). A one-port end
piece uses only its incident length.

O-minimal uniform finiteness gives a constant \(C(\mathrm{base})\)
bounding the number of connected components in intersections of any
member of this local family with any affine line. For a line transverse
to the final smooth hypersurface the intersections are isolated, hence
their number is bounded by the same constant. This includes the two
collars and the fixed base piece in one base ball. The constant does not
depend on N, epsilon, or the lengths of the handles.

Together with C and B,

\[
 i(\Sigma_N)\le 2C(\text{base})+32.
 \tag{E2}
\]

\hypertarget{1022-f.-neumann-spectra-of-the-base-pieces-converge-to-the-original-base}{%
\section{F. Neumann spectra of the base pieces converge to the original
base}\label{1022-f.-neumann-spectra-of-the-base-pieces-converge-to-the-original-base}}

For a fixed finite collection of port locations, remove the one or two
disks of radius \(2\varepsilon\) used by the chain and add the collars
E1. The first and last copies use one port each; all interior copies use
two, and unused ports remain part of the original base. Cut at the tube
ends and impose Neumann conditions there. Denote this compact
manifold-with-boundary by \(B_\varepsilon\).

The standard extension/Rellich argument gives

\[
 \lambda_1^N(B_\epsilon)\longrightarrow\lambda_1(B),\qquad
 \operatorname{Vol}(B_\epsilon)\longrightarrow\operatorname{Vol}(B).
 \tag{F1}
\]

Here is the required compactness mechanism, to make clear that no long
thin region is hidden in this claim. The collars have diameter
\(O(\varepsilon )\) and, after scaling, a uniformly controlled fixed
geometry. Uniform extension from a Euclidean annulus to its ball, chosen
to preserve constants and then transported through local coordinate
charts, extends functions on the punctured base to the whole base with
uniformly bounded \(H^1\) norm. For clarity, write \(\Gamma_\epsilon\)
for the common collar/annulus boundary. The fixed-domain trace and
Poincaré inequalities, after scaling, give

\[
 \int_{\mathrm{collar}} |f|^2
 \le C\epsilon\int_{\Gamma_\epsilon}|f|^2
       +C\epsilon^2\int_{\mathrm{collar}}|\nabla f|^2,
\]

and

\[
 \epsilon\int_{\Gamma_\epsilon}|f|^2
 \le C\int_{\mathrm{annulus}}|f|^2
       +C\epsilon^2\int_{\mathrm{annulus}}|\nabla f|^2.
\]

Combining them gives

\[
 \int_{\rm collar}|f|^2
 \le C\int_{\rm annulus}|f|^2
   +C\epsilon^2\int_{\rm annulus}|\nabla f|^2
   +C\epsilon^2\int_{\rm collar}|\nabla f|^2.
 \tag{F2}
\]

A family with bounded energy and \(L^2\) norm therefore has vanishing
\(L^2\) mass in the shrinking collars: after extension, Rellich
compactness on the fixed closed base makes the annular \(L^2\) masses
vanish. A bounded-energy orthonormal family converges, after
subsequences, to an orthonormal family on the base. For the energy lower
limit, first apply weak lower semicontinuity on the fixed complement of
port balls of radius delta, where the extension agrees with the original
function once epsilon is small. Then let delta decrease to zero. This
gives the required inequality with the original base-piece energies; the
energy introduced inside the artificially filled holes need not tend to
zero. Applying min--max to the limiting orthonormal family gives the
lower eigenvalue limit. Restriction of smooth test functions, with the
obvious uniformly controlled collar extension, gives the upper limit.
This works also when n=2; no assertion that a long cylindrical neck has
a large gap is used.

For each finite N, only finitely many ports occur. Thus epsilon may be
chosen small enough that (F1) holds at every base piece simultaneously.
No uniform rate in N is needed.

\hypertarget{1022-g.-spectral-counting-and-the-contradiction}{%
\section{G. Spectral counting and the
contradiction}\label{1022-g.-spectral-counting-and-the-contradiction}}

Normalize the fixed base B so that \(\lambda _1(B)=1\); choose it with
volume V as large as needed, using the unbounded normalized
first-eigenvalue theorem for smooth embedded hypersurfaces.
Approximation by an algebraic hypersurface can be made C1-small, so this
remains possible with a definable base.

Fix a sufficiently large a in A. For each N, construct the chain and
then choose epsilon sufficiently small so that every base piece has
first positive Neumann eigenvalue at least \(1/2\), volume at least
\(V/2\), and every handle has at most two Neumann eigenvalues below
\(1/2\).

The construction replaces exactly two disk boundaries for every
connecting tube, so the result is a smooth closed hypersurface. The
chain makes it connected, and the placement and clearance estimates make
it embedded. Cut this hypersurface along all base/handle interfaces. Its
Neumann-bracketed disjoint union has at most

\[
 N+2(N-1)=3N-2
\]

eigenvalues strictly below \(1/2\), counting its zero modes. Since
eigenvalues are indexed from zero, min--max gives

\[
 \lambda_{3N-2}(\Sigma_N)\ge\frac12,
 \qquad
 \operatorname{Vol}_n(\Sigma_N)\ge\frac{NV}{2}.
 \tag{G1}
\]

The intersection index is bounded independently of N by E2. If the
proposed additive estimate held, continuity of F on the compact interval
\([0,2C(base)+32]\) would bound its first term by a finite number
\(M_B\). Thus

\[
 \frac12\left(\frac{NV}{2}\right)^{2/n}
 \le M_B+A_n(3N-2)^{2/n}.
\]

Divide by \(N^{2/n}\) and let N tend to infinity:

\[
 \frac12\left(\frac V2\right)^{2/n}\le A_n3^{2/n}.
\]

Since the normalized base volume V can be arbitrarily large, this is
impossible for a dimension-only finite \(A_n\).

The exact numerical constants 2, 8, 32, and 3 are irrelevant; their
independence of N is what permits the contradiction.

\hypertarget{1022-h.-two-strengthenings}{%
\section{H. Two strengthenings}\label{1022-h.-two-strengthenings}}

Continuity of the proposed additive term is not needed for this
contradiction. The intersection index is an integer and stays in the
finite set \(\{0,1,\ldots,\lfloor 2C(\mathrm{base})+32\rfloor\}\). Any
real-valued function on the nonnegative integers is bounded on that
finite set. Thus the same argument rules out an arbitrary finite-valued
additive function of the index.

For \(n\ge 3\), the counterexample can also be confined to hypersurfaces
diffeomorphic to \(S^n\). In Theorem 1.4(2) of {[}1{]}, take the fixed
original submanifold to be the standard embedded sphere. The selected
base is then an n-sphere, and the algebraic approximation preserves its
diffeomorphism type. Joining N such components along a chain is their
connected sum, again diffeomorphic to \(S^n\). The large normalized
first eigenvalue input in dimension two requires varying topology, so
this additional fixed-topology conclusion is asserted only for
\(n\ge 3\).

\hypertarget{1022-references}{%
\section{References}\label{1022-references}}

{[}1{]} B. Colbois, E. B. Dryden, A. El Soufi, \emph{Bounding the
eigenvalues of the Laplace--Beltrami operator on compact submanifolds},
Bulletin of the London Mathematical Society \textbf{42} (2010), 96--108,
Theorem 1.4. DOI: \url{https://doi.org/10.1112/blms/bdp100}. Primary
preprint: \url{https://arxiv.org/abs/0909.5346}.

{[}2{]} L. van den Dries, A. Macintyre, D. Marker, \emph{The elementary
theory of restricted analytic fields with exponentiation}, Annals of
Mathematics \textbf{140} (1994), 183--205. DOI:
\url{https://doi.org/10.2307/2118545}. This establishes o-minimality of
the restricted analytic structure expanded by the real exponential
function.

{[}3{]} L. van den Dries, C. Miller, \emph{On the real exponential field
with restricted analytic functions}, Israel Journal of Mathematics
\textbf{85} (1994), 19--56. The uniform finite-component bound used in
Section E is the standard consequence of definable cell decomposition:
for a fixed definable family, the number of connected components of
every fiber is bounded uniformly over its parameters. A precise statement
is Section~4.4 (PDF p.~13) of [4].


{[}4{]} L. van den Dries and C. Miller,
\emph{Geometric categories and o-minimal structures},
Duke Mathematical Journal \textbf{84} (1996), 497--540;
author's revised manuscript dated 20 February 2001.
Section~4.4 gives uniform bounds on connected components of fibers.
\url{https://people.math.osu.edu/miller.1987/newg.pdf}.
Only this uniform-fiber result is used here.


{[}5{]} L. P\'etiard, \emph{Counter-example to a Kr\"oger type spectral inequality},
arXiv:1706.09189v1 (2017), six-page preprint.
The additive intersection-index question is on p.~3, immediately before
Section~2.
\url{https://arxiv.org/abs/1706.09189v1}.

The polynomial approximation, min--max principle, coarea-free
one-dimensional handle estimates, and elementary extension/Rellich
argument are presented in the proof. No numerical approximation is used.
\end{document}
